\chapter{单变量函数的积分学}
\section{积分}
\subsection{积分的定义}
\begin{definition}
    对于一个函数$f(x),x\in[a,b]$，在$a$与$b$之间插入$n-1$个分割点$x_1,x_2,\cdots,x_{n-1}$，并记$x_0 = a,x_n = b$，称上述过程为区间$[a,b]$的一个\overtext{分割}{partition}，记作：
    \[T:a=x_0<x_1<\cdots<x_{n-1}<x_n=b\]
    这个过程将$[a,b]$分割为$n$个小区间$[x_{i-1},x_i]$，设有：
    \[\left\|T\right\| = \max_{1\leqslant i \leqslant n}\{x_i-x_{i-1}\} = \max_{1\leqslant i \leqslant n}\Delta x_i,\quad i=1,2,\cdots,n\]
    称$\left\|T\right\|$为分割的\overtext{模}{norm}.

    在每个区间$[x_{i-1},x_i]$上任取一点$\xi_i$，设有：
    \[S_n(T) = \sum_{i=1}^{n} f(\xi_i)\Delta x_i\]
    称$S_n(T)$是$f(x)$在$[a,b]$上对应于分割$T$的一个\textbf{Riemann和}.
\end{definition}

\begin{definition}
    若有$I \in \mathbb{R}$使得对$\forall \varepsilon>0$，都$\exists \delta > 0$在$\left\|T\right\|<\delta$时满足：
    \[|S_n(T) - I| < \varepsilon \quad\text{或}\quad \lim\limits_{\left\|T\right\| \to 0} S_n(T) = I\]
    那么称$f(x)$在$[a,b]$上\textbf{Riemann}\overtext{可积}{integrable}，记作：
    \[I = \int_{a}^{b} f(x)\dd{x}\]
    称此定义下的积分为\textbf{Riemann积分}，它是\overtext{定积分}{definite integral}的一种定义.如果不特别注明，「可积」一般指的都是「Riemann可积」.
\end{definition}
上述的定义中包含了两个任意性，一方面是分割取点的任意性，另一方面是取小区间高度的任意性，即$\xi_i \in [x_{i-1},x_i]$的任意性.


\subsection{可积函数类}
\begin{theorem}
    若$\left.f\right|_{[a,b]}$可积，则$\left.f\right|_{[a,b]}$有界.
\end{theorem}

\begin{proof}
    根据Riemann积分的定义，不妨取$\varepsilon = 1$，必然$\exists \delta_1>0$在$\left\|T\right\| < \delta_1$时有：
    \[I-1 < \sum_{i=1}^{n} f(\xi_i)\Delta x_i < I+1\]
    对$\forall \xi_i \in [x_{i-1},x_i]$成立.

    假设$\left.f\right|_{[a,b]}$无界，至少$\exists [x_{k-1},x_k] \subset [a,b]$使得$f$在其上无界.于是将其它有界的部分剥离，有：
    \[I - 1 - \sum_{i \neq k}f(\xi_i)\Delta x_i < f(\xi_k)\Delta x_k < I + 1 - \sum_{i \neq k}f(\xi_i)\Delta x_i\]
    现在取定$\xi_i = x_i$，此处$i = 1,2,\cdots,n$但$i \neq k$.则对于上式，左右两边将是一个固定的数，即：
    \[|f(\xi_k)|\Delta x_k < M\]
    与$\left.f\right|_{[x_{k-1},x_k]}$无界矛盾，故而对于在$[a,b]$上可积的$f$，其在$[a,b]$上必有界.
\end{proof}

接下来列出一些重要的可积函数类，此处暂不证明.
\begin{theorem}
    \label{thm:ClassOfIntegrableFunctions}
    \begin{enumerate}[circ]
        \item 在闭区间$[a,b]$上连续的函数，在$[a,b]$上必可积\textbf{(连续必可积)}.
        \item 若$\left.f\right|_{[a,b]}$有界，且该区间上至多只有有限个间断点，则$\left.f\right|_{[a,b]}$上可积.
        \item 若$\left.f\right|_{[a,b]}$上单调，即使有无限个间断点，$\left.f\right|_{[a,b]}$上也可积\textbf{(单调必可积)}.
    \end{enumerate}
\end{theorem}

在上述诸定理中，\circnum{2}是对\circnum{1}的扩充，这表明可积函数类是大于连续函数类的，可积比连续的宽容度要高.

\begin{extension}{测度论与Lebesgue定理}
    对于Dirichlet函数
    $D(x) = \begin{cases}
        1, \quad x \in \mathbb{Q}\\
        0, \quad x \in \mathbb{R\setminus Q}
    \end{cases}$，其是有界的，但其Riemann不可积，这是因为其处处间断，有无限个间断点.又注意到在\cref{thm:ClassOfIntegrableFunctions}的\circnum{3}中提到，如果$\left.f\right|_{[a,b]}$上单调，即使有无限个间断点也Riemann可积.

    这表明，间断点的「无限多」亦有差异，是否能够通过精确刻画「间断点集合的大小」来判断函数是否Riemann可积？研究这个问题的数学领域称为\overtext{测度论}{measure theory}.

    可以认为\overtext{测度}{measure}是描述集合的「尺度」的一种概念.例如对于集合$E=[a,b]$，其Lebesgue测度$m(E)=b-a$.
    
    在测度的基础上，判断具有间断点函数是否Riemann可积的条件是：
    \begin{theorem}
        \textbf{Lebesgue定理}

        有界函数在$[a,b]$上Riemann可积，当且仅当其间断点的集合测度为$0$.
    \end{theorem}

    在Dirichlet函数中，由于无理数和有理数的稠密性，其处处有间断点，即其在$[0,1]$上间断点集合的测度为$1 \neq 0$，从而其Riemann不可积.

    \begin{definition}
        设有集合$E$，若对$\forall \varepsilon > 0$，都$\exists$可数个开区间$I_n$使得：
        \[E \subset \bigcup_{n=1}^{\infty} I_n \quad\text{且}\quad \sum_{n=1}^{\infty} |I_n| < \varepsilon\]
        称$E$是一个\overtext{零测集}{null set}，其可以被任意小总长度的可数个开区间覆盖.
    \end{definition}
    
    对于零测集，其有一些奇妙的特性，例如对于$E=\{r_1,r_2,\cdots\}$是$[0,1]$上的有理数集，$\forall \varepsilon > 0$，设有：
    \[I_n = \left(r_n - \frac{\varepsilon}{2^{n+1}},r_n + \frac{\varepsilon}{2^{n+1}}\right)\]
    由于有理数是可数的，所以$I_n$是可数的，根据有理数的稠密性显然可以得知$E$是零测集：
    \[E \subset \bigcup_{n=1}^{\infty} I_n \quad\text{且}\quad \sum_{k=1}^{n} |I_k| = \sum_{k=1}^{n} \frac{\varepsilon}{2^n} = \left(1-\frac{1}{2^n}\right)\varepsilon < \varepsilon\]
    这表明即使有理数稠密，把区间中所有有理数及其附近的任意小区间都抽掉，也只是抽走了任意小的长度！
\end{extension}


\subsection{积分的初等例子}
\label{subsec:ExampleOfIntegral}
\begin{question}
    求$f(x) = \sin x$在区间$[a,b]$上的积分.
\end{question}

\begin{solution}
    由于$f(x)$是连续函数，其必然Riemann可积，取分割点：
    \[a,a+h,\cdots,a+(n-1)h,b=a+nh\]
    这使得$[a,b]$被分为$n$等份，有$h=\dfrac{b-a}{n}$.取每个小区间右端点作为$\xi_i$，则Riemann和为：
    \[S_n = h(\sin(a+h) + \sin(a+2h) + \cdots + \sin(a+nh))\]
    对上式中括号内的项乘以$2\sin\dfrac{h}{2}$，利用积化和差公式可以裂项相消：
    \begin{align*}
        \Rightarrow\quad  S_n &= \frac{h}{2\sin\frac{h}{2}}\left(\cos\left(a+\frac{h}{2}\right) - \cos\left(a+\frac{2n+1}{2}h\right)\right)\\
        &= \frac{h}{2\sin\frac{h}{2}}\left(\cos\left(a+\frac{h}{2}\right) - \cos\left(b+\frac{h}{2}\right)\right)
    \end{align*}
    当$h \to 0$时有：
    \[S_n \to (-\cos b) - (-\cos a) = \int_{a}^{b}\sin x\dd{x}\]
\end{solution}

应当注意的是，在求取Riemann和之前，应当先行证明函数Riemann可积.

\begin{question}
    \label{ques:DefinitionOfJ(x)}
    设$a<b$，而$c$满足$a \leqslant c \leqslant b$，考虑如下函数$J(x)$，证明$\displaystyle\int_{a}^{b}J(x)\dd{x} = 0$.
    \[J(x) = \begin{cases}
        0,\quad x \in [a,b],x\neq c\\
        1,\quad x = c
    \end{cases}\]
\end{question}

\begin{proof}
    取$[a,b]$上的任意分割$T:a=x_0<x_1<\cdots<x_n=b$，不妨设$c \in [x_{l-1},x_l]$，对$\forall \xi_i \in [x_{i-1},x_i],1\leqslant i \leqslant n$，有：
    \[\left|\sum_{i=1}^{n}J(\xi_i)\Delta x_i - 0\right| = |J(\xi_l)\Delta x_l - 0| \leqslant \Delta x_l \leqslant \left\|T\right\|\]
    则$\forall \varepsilon > 0$，只需取$\delta = \varepsilon$，在$\left\|T\right\| < \delta = \varepsilon$时有：
    \[\left|\sum_{i=1}^{n}J(\xi_i)\Delta x_i - 0\right| \leqslant \left\|T\right\| < \delta = \varepsilon\]
    对$\forall \xi_i \in [x_{i-1},x_i]$成立，即$\displaystyle\int_{a}^{b}J(x)\dd{x} = 0$.
\end{proof}

上述证明表明，单个点下的面积为$0$.


\subsection{积分的基本性质}
\begin{theorem}
    设有$a<c<b$，若函数$f$在区间$[a,b]$上可积，则$f$在区间$[a,c],[c,b]$上均可积；若$f$在$[a,c],[c,b]$上均可积，则$f$在$[a,b]$上可积，且有：
    \[\int_{a}^{b}f(x)\dd{x} = \int_{a}^{c}f(x)\dd{x} + \int_{c}^{b}f(x)\dd{x}\]
    可以推出，$f$在$[a,b]$的任意子区间$[a',b'] \subseteq [a,b]$上也可积.
\end{theorem}

\begin{theorem}
    设函数$f(x)$和$g(x)$在$[a,b]$上可积，则：
    \begin{itemize}
        \item 对任意常数$\alpha,\beta$，函数$\alpha f(x) + \beta g(x)$在$[a,b]$上可积且有：
        \[\int_{a}^{b}(\alpha f(x) + \beta g(x))\dd{x} = \alpha\int_{a}^{b}f(x)\dd{x} + \beta\int_{a}^{b}g(x)\dd{x}\]
        特别令$\beta \equiv 0$，有：
        \[\int_{a}^{b}\alpha f(x)\dd{x} = \alpha\int_{a}^{b}f(x)\dd{x}\]
        \item 函数$f(x)\cdot g(x)$在$[a,b]$上可积，一般而言：
        \[\int_{a}^{b}f(x)g(x)\dd{x} \neq \int_{a}^{b}f(x)\dd{x}\cdot\int_{a}^{b}g(x)\dd{x}\]
    \end{itemize}
\end{theorem}

在上述定理的基础上，引入\cref{subsec:ExampleOfIntegral}中的\cref{ques:DefinitionOfJ(x)}中所定义的$J(x)$，注意到：
\[\int_{a}^{b}\alpha J(x)\dd{x} = 0 \quad\Rightarrow\quad \int_{a}^{b}(f(x) + \alpha J(x))\dd{x} = \int_{a}^{b}f(x)\dd{x}\]
这样的结果表明：
\begin{theorem}
    改变可积函数$f(x)$有限个点处的值，不会破坏其可积性，也不会改变其积分值.
\end{theorem}

\begin{theorem}
    设函数$f(x)$和$g(x)$在$[a,b]$上可积：
    \begin{itemize}
        \item \textbf{积分的保号性：}若对$\forall x \in [a,b]$有$f(x) \geqslant 0$，则$\displaystyle\int_{a}^{b}f(x)\dd{x}\geqslant 0$.
        \item \textbf{积分的单调性：}若对$\forall x \in [a,b]$有$f(x)\leqslant g(x)$，则$\displaystyle\int_{a}^{b}f(x)\dd{x} \leqslant \int_{a}^{b}g(x)\dd{x}$.
        \item 函数$|f(x)|$在$[a,b]$上也可积且$\displaystyle\left|\int_{a}^{b}f(x)\dd{x}\right|\leqslant\int_{a}^{b}|f(x)|\dd{x}$.
    \end{itemize}
\end{theorem}

上述诸定理将在后续得到证明.

\begin{theorem}
    \textbf{积分中值定理}

    设函数$\left.f\right|_{[a,b]}$连续，则$\exists \xi \in [a,b]$使得：
    \[\int_{a}^{b}f(x)\dd{x} = f(\xi)(b-a)\]
\end{theorem}

\begin{proof}
    由于$\left.f\right|_{[a,b]}$连续，则其在$[a,b]$上必有最大最小值：
    \[m \leqslant f(x) \leqslant M\]
    两边积分，有：
    \[m \leqslant \frac{1}{b-a}\int_{a}^{b}f(x)\dd{x} \leqslant M\]
    根据连续函数的介值定理，对于介于$m$与$M$之间的一个值$\displaystyle\frac{1}{b-a}\int_{a}^{b}f(x)\dd{x}$，必有$\xi \in [a,b]$使得$f(\xi)$等于这个值.
\end{proof}


\subsection{微积分基本定理}
\begin{definition}
    设$\left.f\right|_{[a,b]}$可积，则对于$a\leqslant x \leqslant b$有$\left.f\right|_{[a,x]}$也可积，记：
    \[\varphi: x \mapsto \int_{a}^{x}f(t)\dd{t} \quad\text{或}\quad \varphi(x) = \int_{a}^{x}f(t)\dd{t}\]
    称$\varphi(x)$为$f$的\overtext{变上限积分}{integral with variable upper limit}，其中对应$f$的自变量可以随意变换而不影响外部变量，称之为\overtext{哑变量}{dummy variable}.
\end{definition}

\begin{theorem}
    设$\left.f\right|_{[a,b]}$可积，则$f$的变上限积分$\varphi(x)$在$[a,b]$上连续.
\end{theorem}

\begin{proof}
    由于可积必有界，因此$\exists M>0$使得$|f(t)|\leqslant M,t\in [a,b]$.设有$x_0\in[a,b]$，$\forall x \in[a,b]$有：
    \begin{align*}
        |\varphi(x) - \varphi(x_0)| &= \left|\int_{a}^{x}f(t)\dd{t} - \int_{a}^{x_0}f(t)\dd{t}\right|\\
        &= \left|\int_{x_0}^{x}f(t)\dd{t}\right| \leqslant \left|\int_{x_0}^{x}|f(t)|\dd{t}\right| \leqslant \left|\int_{x_0}^{x}M\dd{t}\right| = M|x-x_0|
    \end{align*}
    要使上式在$|x-x_0|<\delta$时小于$\forall \varepsilon > 0$，只需取$\delta = \dfrac{\varepsilon}{M}$，就有$|\varphi(x)-\varphi(x_0)|<\varepsilon$，进而证明了$\varphi(x)$在$[a,b]$上连续.
\end{proof}

\begin{theorem}
    \label{thm:VariableLimitIntegralDerivative}
    设函数$\left.f(x)\right|_{[a,b]}$上可积，设$x_0 \in [a,b]$：
    \begin{itemize}
        \item 若$f(x)$在$x=x_0$处连续，则$f(x)$在$x_0$处的变上限积分$\varphi(x)$在$x=x_0$处可微且$\varphi'(x_0) = f(x_0)$.
        \item 若$f(x)$在$[a,b]$上处处连续，则$f(x)$在$[a,b]$上的变上限积分$\varphi(x)$在$[a,b]$上可微且$\varphi'(x) = f(x)$，即连续函数一定有原函数：
        \[F(x) = \int f(x)\dd{x} = \varphi(x) + C = \int_{a}^{x} f(t)\dd{t} + C\]
    \end{itemize}
\end{theorem}

\begin{proof}
    要证$\varphi(x)$在$x=x_0$可微，只需证存在$\lim\limits_{x \to x_0}\dfrac{\varphi(x)-\varphi(x_0)}{x-x_0} = f(x_0)$.
    \begin{align*}
        \left|\frac{\varphi(x) - \varphi(x_0)}{x-x_0} - f(x_0)\right| &= \left|\frac{1}{x-x_0}\left(\int_{a}^{x}f(t)\dd{t} - \int_{a}^{x_0}f(t)\dd{t} - \int_{x_0}^{x}f(x_0)\dd{t}\right)\right|\\
        &= \left|\frac{1}{x-x_0}\int_{x_0}^{x}(f(t)-f(x_0))\dd{t}\right|\\
        &\leqslant \frac{1}{|x-x_0|}\left|\int_{x_0}^{x}(f(t)-f(x_0))\dd{t}\right|
    \end{align*}
    由于$f(x)$在$x=x_0$处连续，故而$\forall \varepsilon > 0$，$\exists \delta > 0$在$|x-x_0|<\delta$时有$|f(x)-f(x_0)|<\varepsilon$，则有：
    \begin{align*}
        \frac{1}{|x-x_0|}\left|\int_{x_0}^{x}(f(t)-f(x_0))\dd{t}\right| &\leqslant \frac{1}{|x-x_0|}\left|\int_{x_0}^{x}|f(t)-f(x_0)|\dd{t}\right|\\
        &< \frac{1}{|x-x_0|}\left|\int_{x_0}^{x}\varepsilon\dd{t}\right| \leqslant \varepsilon
    \end{align*}
    故而得证$\lim\limits_{x \to x_0}\dfrac{\varphi(x)-\varphi(x_0)}{x-x_0} = f(x_0)$.对于在$[a,b]$上处处连续的情况，将上述证明推广即得证.
\end{proof}

在上述证明的基础上，最终可以给出：

\begin{theorem}
    \textbf{Newton-Leibniz公式(微积分基本定理)}

    设$f(x)$是$[a,b]$上的连续函数，若$F(x)$是$f(x)$的一个原函数，则：
    \[\int_{a}^{b} f(x)\dd{x} = \int_{a}^{b}\dd{F(x)} = F(b) - F(a) = F(x)\Big|_{a}^{b}\]
\end{theorem}

需要注意，上述Newton-Leibniz公式的证明基于\cref{thm:VariableLimitIntegralDerivative}中所假设的$f(x)$在$[a,b]$上处处连续，实际上，利用Riemann和，可以给出更弱化条件下的Newton-Leibniz公式：

\begin{theorem}
    设$\left.f(x)\right|_{[a,b]}$上可积，若$f(x)$有原函数$F(x)$，则满足：
    \[\int_{a}^{b}f(x)\dd{x} = F(b) - F(a)\]
\end{theorem}

\begin{proof}
    设区间$[a,b]$上有任意的分割：
    \[T:a=x_0<x_1<\cdots<x_{n-1}<x_n=b\]
    则有：
    \[\sum_{i=1}^{n}(F(x_i) - F(x_{i-1})) = F(b) - F(a)\]
    由于$F(x)$作为原函数，其在$[a,b]$上处处连续且可导，其在每个小区间$[x_{i-1},x_i]$上满足Lagrange中值定理，因此必然$\exists \xi_i \in (x_{i-1},x_i)$使得：
    \[F(x_i) - F(x_{i-1}) = F'(\xi_i)\Delta x_i = f(\xi_i)\Delta x_i \quad\Rightarrow\quad \sum_{i=1}^{n}f(\xi_i)\Delta x_i = F(b)-F(a)\]
    根据Riemann积分的定义，有：
    \[\int_{a}^{b}f(x)\dd{x} = \lim_{\left\|T\right\| \to 0}\sum_{i=1}^{n}f(\xi_i)\Delta x_i = F(b)-F(a)\]
\end{proof}

但是还存在一些例外，并非所有函数都满足Newton-Leibniz公式的条件：
\begin{enumerate}[circ]
    \item 有的函数可积却没有原函数，如\cref{subsec:ExampleOfIntegral}\cref{ques:DefinitionOfJ(x)}中的$J(x)$.
    \item 有的函数因其连续而有原函数，却不能用有限次初等函数的组合表示，如$f(x) = \ee^{-x^2}$.
\end{enumerate}
对于情况\circnum{1}，无法使用Newton-Leibniz公式求解，通常需要基于Riemann积分的定义来求解；而对于情况\circnum{2}，需要引入一些特殊函数或积分来表示其原函数，再使用Newton-Leibniz公式求解.

这表明，Newton-Leibniz公式成立的基本条件是存在原函数，而不是仅仅要求函数可积.